% ============================================================
% Sección 1
% ============================================================
\section{1 Métodos Numéricos en \texttt{solve\_ivp} de Python}

\begin{tcolorbox}[enunciado]
Está muy extendido en el imaginario colectivo de los estudiantes que las computadoras lo hacen todo bien y que basta con poner el código para que en automático todo salga perfecto, no obstante ésta percepción es incorrecta como de manera muy sencilla veremos en esta sección.

Además se te hacen algunas preguntas, que debes contestar siguiendo la bibliografía indicada.

\vspace{0.2cm}
\textbf{Clases (Apoyo Secundario):}
\begin{itemize}[leftmargin=*]
    \item Solución EDO con y sin unicidad 24Ago26 MCA4
    \item Sol Equilibrio, Curva Integral, Campo Direccional 26 Ago 2026 MCA4
\end{itemize}

\textbf{Código:}
\begin{itemize}[leftmargin=*]
    \item Github: \texttt{1.1\_Campo\_Direccional\_y\_Curva\_Integral}
    \item GiHub: \texttt{2.3\_Métodos\_Numericos}
\end{itemize}

\textbf{Documentación:}
\begin{itemize}[leftmargin=*]
    \item Web: Python Doc \texttt{solve\_ivp}
\end{itemize}

\textbf{Requisitos:} \\
Debes instalar numba:
\begin{terminal}
pip install numba
\end{terminal}
\end{tcolorbox}

\subsection*{ACTIVIDADES}

\begin{enumerate}[label=\arabic*., leftmargin=*]
    \item En las clases arriba mencionadas y en la Actividad 1 sección 3.3 vimos la EDO:
    \[ \frac{dy}{dx} = \frac{4-2x}{3y^{2}-5} \]
    cuya función (ecuación) solución implícita es:
    \[ y^{3}-5y+x^{2}-4x=c \]
    y en el código 
    \[\texttt{Campo\_Direccional\_Curvas\_Integrales\_dy\_4-2x\_div\_3y2-5.py} \]
    se grafican el campo direccional y un par de curvas integrales (solución) de la EDO para distintos valores de $c$.

    Usando la documentación oficial responde y realiza lo que se indica en los siguientes incisos:

    \begin{enumerate}[label=(\alph*), leftmargin=*]
        \item Transcriba y ejecute el código.
        \subsection*{Solución}
        % Escribe tu respuesta aquí

        \item ¿Cuáles son los métodos que puede utilizar \texttt{solve\_ivp}?
        \subsection*{Solución}
        % Escribe tu respuesta aquí

        \item Si no se indica algún método particular, ¿cuál se utiliza?
        \subsection*{Solución}
        % Escribe tu respuesta aquí

        \item En la solución: \texttt{sol1}
        \begin{enumerate}[label=\roman*., leftmargin=*]
            \item Cambia \texttt{BDF} por \texttt{RK45}, es decir, modifica:
            \begin{terminal}
sol1 = solve\_ivp(f, (0,4), [1.3], method='BDF')
            \end{terminal}
            por:
            \begin{terminal}
sol1 = solve\_ivp(f, (0,4), [1.3], method='RK45')
            \end{terminal}
            ¿Con cuál método consideras se calcula adecuadamente la solución?
            \subsection*{Solución}
            % Escribe tu respuesta aquí

            \item Analiza el número de pasos tomado por \texttt{BDF} y \texttt{RK45}. \\
            \textbf{HINT:} Puedes ver el ``explorador de variables'' o utilizar la instrucción \texttt{len(sol1.t)} (muy recomendable). ¿Qué método tomó más pasos?
            \subsection*{Solución}
            % Escribe tu respuesta aquí

            \item Analiza la distribución de puntos calculados por método de manera gráfica. Para esto solo debes graficar los puntos calculados. Cambia:
            \begin{terminal}
plt.plot(sol1.t, sol1.y[0], color='brown')
            \end{terminal}
            a la forma:
            \begin{terminal}
plt.plot(sol1.t, sol1.y[0], 'o', color='brown', markersize=3)
            \end{terminal}
            Usando los resultados obtenidos, justifica por qué un método lo hace bien y el otro incorrectamente.
            \subsection*{Solución}
            % Escribe tu respuesta aquí
        \end{enumerate}

        \item ¿Qué método utiliza \texttt{solve\_ivp} para \texttt{sol3}?
        \subsection*{Solución}
        % Escribe tu respuesta aquí

        \item En solución: \texttt{sol4}
        \begin{enumerate}[label=\roman*., leftmargin=*]
            \item Cambia \texttt{LSODA} por \texttt{DOP853}, es decir, modifica:
            \begin{terminal}
sol4 = solve\_ivp(f, (-0.8, 4.6), [-2.6], method='LSODA')
            \end{terminal}
            por:
            \begin{terminal}
sol4 = solve\_ivp(f, (-0.8, 4.6), [-2.6], method='DOP853')
            \end{terminal}
            ¿Qué modificación presenta la gráfica de la curva solución?
            \subsection*{Solución}
            % Escribe tu respuesta aquí

            \item Por simple inspección, ¿cómo justificas que el resultado dado por el método ``\texttt{DOP853}'' no es correcto?
            \subsection*{Solución}
            % Escribe tu respuesta aquí

            \item Analiza el número de pasos tomado por \texttt{LSODA} y \texttt{DOP853}. \\
            \textbf{HINT:} Puedes utilizar la instrucción \texttt{len(sol4.y[0])}. ¿Cuántos pasos genera cada método?
            \subsection*{Solución}
            % Escribe tu respuesta aquí
        \end{enumerate}

        \item De manera breve y siguiendo la documentación base, di a qué se refieren cada uno de los 6 métodos que implementa \texttt{solve\_ivp}.
        \subsection*{Solución}
        % Escribe tu respuesta aquí
    \end{enumerate}

    \item Reescribe y ejecuta el código, dado en Github Métodos numéricos: 
    \[\texttt{RK4\_Oscilador\_Armonico\_Py\_vs\_Numba.py}\]
    \begin{enumerate}[label=(\alph*), leftmargin=*]
        \item ¿Cuál es el tiempo de solución al problema con:
        \begin{itemize}[leftmargin=*]
            \item \texttt{RK4 Python}
            \item \texttt{RK4 Numba}
        \end{itemize}
        \textbf{HINT:} Esto lo imprime el código.
        \subsection*{Solución}
        % Escribe tu respuesta aquí

        \item ¿Cuál es la diferencia de tiempo entre ambos procedimientos?
        \subsection*{Solución}
        % Escribe tu respuesta aquí

        \item ¿El problema que resuelven es el mismo? Menciona cuál es.
        \subsection*{Solución}
        % Escribe tu respuesta aquí

        \item Los métodos \texttt{rk4\_python} y \texttt{rk4\_numba} ¿son iguales?, en el sentido de que:
        \begin{itemize}[leftmargin=*]
            \item reciben los mismos parámetros
            \item la estructura interna (codificación) es igual
            \item se le manda llamar con los mismos valores
        \end{itemize}
        \subsection*{Solución}
        % Escribe tu respuesta aquí

        \item Habiendo resuelto lo anterior, justifica por qué un procedimiento es más rápido (eficiente) que el otro. \\
        \textbf{HINT:} Qué hace \texttt{numba} y la instrucción:
        \begin{terminal}
@njit(cache=True, fastmath=True)
        \end{terminal}
        \subsection*{Solución}
        % Escribe tu respuesta aquí
    \end{enumerate}
\end{enumerate}